% LSPV-2 REPORT::CIRCUIT-PRIMARY-SOURCE-SELECT174::2026-10-27T05:12:00Z
\documentclass[11pt]{article}
\usepackage[T1]{fontenc}
\usepackage[utf8]{inputenc}
\usepackage[margin=.82in]{geometry}
\usepackage{amsmath,amssymb,amsthm,booktabs,longtable,xurl,hyperref}
\hypersetup{hidelinks}
\newtheorem{theorem}{Theorem}
\newtheorem{lemma}[theorem]{Lemma}
\newtheorem{counterexample}[theorem]{Counterexample}
\newcommand{\SAT}{\mathsf{SAT}}
\newcommand{\OPEN}{\textbf{OPEN}}
\newcommand{\Select}{\mathsf{Select}}
\newcommand{\Rank}{\mathsf{Rank}}
\title{CIRCUIT--SCAN--REVERSE--SELECT173\\
\large Exact, Uniform Inverse Consumers at the Last Forward Comparator,\\
Scan Exit and Reverse-Payload Output of an Oblivious SAT Witness Verifier}
\author{LSPV-2: Circuit / Reduction Builder}
\date{27 October 2026, 04:12 UTC}
\begin{document}\maketitle\n\begin{center}\textbf{Run 174 extended report: source aliases, compiler clones, padding and pin source inverses.}\end{center}\n
\begin{center}
\fbox{\parbox{.94\linewidth}{\textbf{PROVED/EXACT, architecture-restricted; independent referee pending.}
The last forward layer's record-output gates, the scan-result gates, and
all reverse-payload MUX output gates have explicit table-free, signed
direct-consumer and proper-source occurrence Rank/Select with
$O(\log^3 G)$ bit time and $O(\log G)$ combined writable/index bits.
\textbf{COMPUTATIONAL:} 6 complete circuit instances, 1,272 target gates,
3,058 direct uses, 18,592 signed source occurrences, all 8,192
length-10 verifier input/witness pairs, and an eight-clause UNSAT
control, with zero assertion failures.
\textbf{OPEN:} arbitrary source alias/constant/other gate classes,
complete regularizer inverse, all ordinary SAT encodings, ARC-EDGE-90,
and $P\ne NP$.}}\end{center}

\section{Registered falsifiers, model and evidence order}
Before consulting the primary papers and latest project programs, the
following four mechanisms were specified for falsification:
\begin{description}
\item[] \textbf{H1.} Last-forward record outputs enter only a known
scan block and its immediate successor's key-equality check. All
signed consumers should be an arithmetic function of their gate ID.
\emph{Established here for the declared network.}
\item[] \textbf{H2.} Every scan output and every reverse-payload MUX
output should have a constant-fanout inverse consumer map, including
its final row-checker use and unconsumed positions.
\emph{Established with the terminal zero-fanout correction.}
\item[] \textbf{H3.} A reverse-stage upper output necessarily enters
the next reverse stage through its upper input, and every final
reverse output has two consumers. \emph{Both assertions are false.}
\item[] \textbf{H4.} Table-free access to these verifier source gates
produces a fixed-power smaller ordinary-SAT child. \emph{Unsupported
and false for the earlier explicitly pinned two-level architecture
with its known linear output-size floor.}
\end{description}
The audit then read the original 166 gate emitter, the 168/169 inverse
gate cases, the 171/172 model and source selectors, roles 1/3/4/5,
the three model governance directives and original FLMV,
Pippenger--Fischer, Hennie--Stearns and Williams sources. This report
isolates \emph{new sectors} rather than relabeling the previous
nonfinal-layer and comparator-internal theorems.

\paragraph{Endpoint discipline.}
A claimed achievable $(t,s)$ in the FLMV model means \emph{one fixed}
uniform deterministic total ordinary-SAT decider for every $N$-bit
input, using $N^{t+o(1)}$ time and $N^{s+o(1)}$ writable/index
space. The fixed verifier here has a second witness input, so its
size/access results are not a SAT decision-circuit compression.
$P=NP$ supplies only an unknown polynomial exponent $K$, not the
separate hypothetical seed $d=901/500$ and subpolynomial work space.

\section{Exact circuit, encoding and gate-index arithmetic}
Fix $q\ge3$, $m\ge1$, let $w=\lceil\log_2(q+1)\rceil$ and
$N=m[1+3(w+1)]$. A record contains one active flag and three pairs
of a $w$-bit name and polarity. Inactive clauses are true; an
active clause is false if any name is outside $[1,q]$ or names
repeat; otherwise it is a proper signed ternary clause. The total
witness verifier is
\[
V_{q,m}(x,z)=\bigwedge_{c<m}\left[\neg a_c\vee
  \left(\operatorname{Valid}(c)\wedge
    \bigvee_{j=0}^2[z_{v_{cj}}=s_{cj}]\right)\right].\tag{1}
\]
There is no malformed-input promise for \emph{this} fixed-record
language. An arbitrary variable-width DIMACS or binary ordinary-SAT
parser is a separate, unresolved conversion.

Set $K=2^{\lceil\log_2(q+3m)\rceil}$, $h=\log_2K$,
$J=h(h+1)/2$, $D=w+3$, $H=7(w+2)$, $F=H+6D=13w+32$.
Forward bitonic sorting uses $J$ layers, each with $K/2$
comparators of $F$ Boolean AND/OR gates. Every comparison saves its
switch gate and then creates two $D$-wire three-gate MUX outputs.
The successive global gate counts and boundaries are
\begin{align}
G_f&=J\frac K2 F,\\
G_s&=19+(K-1)(4w+18),\\
G_r&=3KJ,\\
G_c&=m(36w+15)+(m-1),\\
G&=G_f+G_s+G_r+G_c,\qquad
B_c=G_f+G_s+G_r.                              \tag{2}
\end{align}
All gate IDs in what follows are \emph{zero-based} indices in one
fixed, fully specified circuit, not machine time or wire count.

The scan works from sorted record position $p=0$ to $K-1$.
Its global starting gate is
\[
b_p=\begin{cases}G_f,&p=0,\\
G_f+19+(p-1)(4w+18),&p>0.
\end{cases}                                                   \tag{3}
\]
It computes 
$\operatorname{isassign}=\operatorname{EqBits}(\mathrm{kind},0)$,
$\operatorname{isquery}=\operatorname{EqBits}(\mathrm{kind},1)$,
$\operatorname{same}=\operatorname{EqBits}(\mathrm{id}_p,
\mathrm{id}_{p-1})$ for $p>0$,
$\mathrm{keep}=\neg\operatorname{isassign}\wedge
  (\operatorname{same}\wedge\mathrm{carry}_{p-1})$,
$\mathrm{carry}_p=(\operatorname{isassign}\wedge\mathrm{payload}_p)
\vee\mathrm{keep}$, and
$\mathrm{scanout}_p=\operatorname{isquery}\wedge\mathrm{carry}_p$.
The source has exactly 19 gates at $p=0$ and $4w+18$ at every
later $p$. In particular the scan output is an \emph{AND} gate
at
\[
                 s_p=G_f+18+p(4w+18).                       \tag{4}
\]

\section{New exact signed inverse for last-forward outputs}
Let $g_{p,c}$ denote the last sorting layer's OR output at
record position $p$ and field $c\in[0,D)$. For every
$0\le p<K$, it exists and can be located by taking the last
layer's stride $s=1$, low position $i=p\,\&\,\neg1$,
pair rank $r=i/2$, output side $\delta=p\bmod2$, and
\[
  g_{p,c}=\left(J-1\right)\frac K2F+rF+H+3(\delta D+c)+2.
                                                               \tag{5}
\]
A signed use is a triple $(\text{consumer gate},
\text{operand slot},\text{input polarity})$, with operand slots
$0,1$ and polarities $+1,-1$.
\begin{theorem}[Last-forward scan inverse, \textbf{PROVED/EXACT restricted}]
For every such $p,c$, the \emph{complete} direct-consumer
list is the sorted union of the following entries:
\begin{description}
\item[] \emph{Identifier bit} $0\le c<w$:
if $p>0$, use
$(b_p+14+3c,0,+)$ and $(b_p+15+3c,0,-)$; if
$p<K-1$, also use
$(b_{p+1}+14+3c,1,-)$ and
$(b_{p+1}+15+3c,1,+)$.
\item[] \emph{Kind bit} $c=w+t$, $t\in\{0,1\}$:
use $(b_p+3t,0,+)$, $(b_p+3t+1,0,-)$,
$(b_p+7+3t,0,+)$ and $(b_p+8+3t,0,-)$.
\item[] \emph{Payload bit} $c=w+2$:
use $(b_p+d_p,1,+)$, where
$d_0=16$ and $d_p=4w+15$ for $p>0$.
\end{description}
The direct fanout equals $2$ for endpoint identifier fields,
$4$ for interior identifiers and all kind fields, and $1$
for payload fields. These signed lists are uniformly computable
in $O(\log^3G)$ elementary bit operations and $O(\log G)$
total index/writable bits without storing the circuit.
\end{theorem}
\begin{proof}
The verifier's \texttt{EqBits} routine implements each literal
comparison $\neg\operatorname{XOR}(u,v)$ by
$\operatorname{AND}(u,\neg v)$,
$\operatorname{AND}(\neg u,v)$ and their OR.
A sorted identifier is read once by the current position's
same-key computation (as the left argument), except at $p=0$,
and once by the next position (as the right argument), except
at $p=K-1$. Each comparison yields exactly the two indicated
signed uses. Each kind bit participates in the assignment and
query equality tests, giving four signed uses. Every payload
enters exactly one AND with the assignment flag. None of these
record fields is used in reverse routing, which routes only the
new scan-result bit. Equation (3) locates every use.
\end{proof}

\section{New scan-output and reverse-MUX inversion}
For sorting layer $\ell\in[0,J)$, let its size/stride be
$z_\ell=2^{a_\ell}$, $s_\ell=2^{b_\ell}$, enumerated by
$a=1,\ldots,h$ and $b=a-1,\ldots,0$. At position $p$,
set
\[
 i_\ell=p\,\&\,\neg s_\ell,\quad
 r_\ell=s_\ell\left\lfloor\frac{i_\ell}{2s_\ell}\right\rfloor
       +(i_\ell\bmod s_\ell),\quad
 \epsilon_\ell=[p\,\&\,s_\ell\ne0].                      \tag{6}
\]
The \emph{reverse} stage corresponding to forward layer
$\ell$ begins its pair $r$ at
\[
 R_{\ell,r}=G_f+G_s+6(K/2)(J-1-\ell)+6r.                \tag{7}
\]
Write $\mathcal U_\ell(p)$ for the two signed gate-input
uses
\[
 \mathcal U_\ell(p)=
 \begin{cases}
 \{(R_{\ell,r_\ell},1,+),(R_{\ell,r_\ell}+4,1,+)\},
               &\epsilon_\ell=0,\\
 \{(R_{\ell,r_\ell}+1,1,+),(R_{\ell,r_\ell}+3,1,+)\},
               &\epsilon_\ell=1.
 \end{cases}                                                 \tag{8}
\]
The reverse layer's OR result at position $p$ is
$g^{\rm rev}_{\ell,p}=R_{\ell,r_\ell}+3\epsilon_\ell+2$.
\begin{theorem}[Scan/reverse exit Select, \textbf{PROVED/EXACT restricted}]
The complete direct-consumer list of the scan result
$s_p$ is $\mathcal U_{J-1}(p)$. For a reverse result
$g^{\rm rev}_{\ell,p}$, if $\ell>0$ it is
$\mathcal U_{\ell-1}(p)$. If $\ell=0$ and
$p\notin[q,q+3m)$, the list is empty. Finally, if
$\ell=0$ and $p=q+3c+j$, $0\le c<m$, $j\in\{0,1,2\}$,
then its entire signed consumer list is
\[
\{(B_c+c(36w+15)+j(8w+3)+8w,1,-),
\ (B_c+c(36w+15)+j(8w+3)+8w+1,1,+)\}.               \tag{9}
\]
The scan results and nonterminal reverse outputs have
fanout two. The terminal reverse outputs have fanout two
\emph{only} in the original query range, otherwise zero.
All of these selectors are uniformly computable with the same
$O(\log^3G)$ bit time and $O(\log G)$ live bits.
\end{theorem}
\begin{proof}
Every reversed comparison uses precisely two three-gate
MUXs with the \emph{previously saved} switch wire. Its lower
input participates in the first MUX's left AND and second
MUX's right AND, at local gates 0 and 4; its upper input
uses gates 1 and 3. Reverse comparator inputs at a stage
are the output payload positions of the immediately later
reverse stage. Their receiving side must be recomputed from
the \emph{receiving} stride by (6), proving (8), with no
uncharged swap table. After layer zero there is no further
MUX: only the original $3m$ literal-request positions are
read as the second operand of the row-checker's
$\operatorname{Eq}(\text{sign},\text{lookup})$ operation,
whose first two ANDs use the lookup with negative and
positive polarity. The $8w$ offset follows from $w-1$
nonzero gates, $7w$ range-comparison gates and one valid
AND gate within each request block. Assignment and padding
positions are never referenced by the row checker.
\end{proof}

\begin{counterexample}[\textbf{FALSIFIED: side and fanout conservation}]
For $q=3,m=1$, $K=8,J=6$. After reverse layer 5, position
$p=1$ is the \emph{upper} output of a stride-1 pair,
yet enters reverse layer 4 through a \emph{lower} stride-2
input: its actual signed uses are $(1623,1,+)$ and
$(1627,1,+)$, not the same-side pair's gates.
At reverse layer 0, position $p=0$ is an assignment slot.
Its output gate $1715$ has \emph{zero} downstream gate
consumers, and exactly five own-block source-CNF occurrences;
query position $p=3$ has output gate $1724$, two signed gate
uses and thirteen source-CNF occurrences. Thus neither
``old side equals next side'' nor ``all terminal reverse outputs
have degree two'' is correct.
\end{counterexample}

\section{Uniform gate ID and source occurrence inverse}
The input to the accessor is a \emph{global gate ID}, not a
sorted-array location supplied by an oracle. A parameter-only
program reconstructs $G_f,G_s,G_r,G_c$, checks whether the ID
is a last-layer output gate (local offset $H+3u+2$), a scan
output gate (formula (4)), or a reverse output (offset 2 or 5
within its six-gate pair). Layer parameters are recovered by
streaming the $J=O(\log^2K)$ bitonic layer counters, without
allocating an array of layers, gates or occurrences. Gate/output
inverse arithmetic then recovers $p,c,\ell,r$ and the preceding
theorems supply at most four signed direct uses.

Every Boolean gate $g$ is compiled to thirteen proper ternary
clauses with two fresh argument clones. Its physical output
variable is
\[
                 v_g=(N+q+3)+9g+2.                         \tag{10}
\]
For an OR gate $g$, its own occurrences are
$(13g+8,2,-)$ and $(13g+r,1,+)$ for $r=9,10,11,12$;
for an AND gate the displayed signs are reversed.
For each signed direct consumer $(a,k,\tau)$ append precisely
\[
  (13a+4k+j,1,\tau)\quad(j=0,1),\qquad
  (13a+4k+j,1,-\tau)\quad(j=2,3).                           \tag{11}
\]
Sort this constant-size list by physical clause/slot.
The selector's rank is its position in that list. Thus its
signed degree is precisely $5+4f$, where $f$ is gate fanout.
The possible degrees in the newly covered sectors are
$5,9,13,21$. A uniform source-variable-to-gate inverse is
available by division by nine after subtracting the offset
in (10); foreign source-variable classes are not claimed.

\begin{theorem}[Source-CNF Rank/Select, \textbf{PROVED/EXACT restricted}]
For every supported gate output $v_g$, the sequence (10)--(11)
is exactly its physical signed occurrence list in the proper
13-clause-per-gate source compiler, including zero-fanout
terminal outputs. The procedures
$\Select(v_g,r)$ and $\Rank(v_g,c,k)$ take
$O(\log^3G)$ elementary bit operations and
$O(\log G)$ combined live work/index bits. No stored
occurrence table, full gate transcript, or free source reader
is required to \emph{locate} these occurrences.
\end{theorem}
\begin{proof}
The cloned input equivalences contain each signed wire
reference exactly twice positively and twice negatively,
with original-use polarity distinguished in the four physical
rows (11). The gate relation contains its output exactly five
times as shown; fresh clone and padding IDs prevent other
collisions. The direct-consumer theorems are complete for the
declared sectors. A constant-length list and gate-class decoder
therefore suffice. Each of $O(\log^2K)$ streamed iterations
uses schoolbook $O(\log^2G)$ or cheaper arithmetic on
$O(\log G)$-bit addresses, so $O(\log^3G)$ bit time is a
conservative uniform bound. Only constantly many addresses,
indices, and signs are live at once. Computing or emitting a
full $\Theta(G)$-gate circuit is not part of this accessor.
\end{proof}

For audit of \emph{physical bit positions}, one canonical fixed-width
compiler serialization has
$n_0=9G+3N+q+8$ variables,
$m_0=13G+4N+12$ clauses and
$W=\lceil\log_2(\max(n_0,m_0)+1)\rceil$ literal-index bits.
Its self-delimiting header has an eight-bit magic tag, unary
length of the width field, its binary width, and two $W$-bit
counts. If its length is $H_0$, occurrence $(c,k)$ starts at
\[
                 H_0+(3c+k)(W+1),                           \tag{12}
\]
with sign first and $W$-bit literal ID next. This arithmetic
produces the exact stored bit offsets, but a \emph{general}
source clause literal reader still requires its own gate
locator and the unresolved pin/alias/constant interfaces.
The source-input pin group depends on $x$, and it is not
silently treated as a constant or absent.

\section{Finite evidence and falsifying controls}
The standalone Python program independently emits the \emph{whole}
fixed witness-verifier DAG, constructs a table of genuine direct
uses and physical 13-row source occurrences for audit \emph{only},
then compares every selected gate to a separate parameter-only
\texttt{Shape(q,m)} interface with no stored netlist, layer table
or occurrence table. All signed global gate, operand, clause,
literal-position and physical header/field predictions match.
The attempted table-free predictor does not inspect the audit
maps. $q=3,m=1$ is also checked against every one of its $2^{10}$
source strings and every one of eight witnesses, including
inactive clauses, invalid IDs, duplicate IDs, SAT and UNSAT
malformed records. The proper eight-clause ternary blocking CNF
on three variables is independently verified UNSAT against all
eight witnesses, and the compiled circuit rejects them.

\begin{center}\small
\begin{tabular}{rrrrrrrr}
\toprule
$q$ & $m$ & $K$ & total $G$ & forward & scan & reverse & signed occurrences\\\midrule
3&1&8&1824&40&8&48&1400\\
4&1&8&2200&48&8&48&1552\\
3&2&16&5704&80&16&160&3664\\
5&2&16&6876&96&16&160&3984\\
7&2&16&6876&96&16&160&3984\\
6&3&16&7000&96&16&160&4008\\\midrule
\multicolumn{4}{r}{\textbf{TOTAL}}&456&80&736&\textbf{18,592}\\\bottomrule
\end{tabular}
\end{center}
The total is 1,272 newly covered Boolean output gates with
3,058 true signed direct gate-input uses. There were 18,592
correct signed source occurrence selectors and inverse-rank
roundtrips, with matching sign and every index bit of each
corresponding fixed-width stored literal. No failed assertions.
The empirical counts are finite testing, not asymptotic SAT
lower bounds or measured Turing-machine transitions.

\section{Three models, same-machine accounting, and ARC firewalls}
\paragraph{RAT or bounded-word RAM.}
For a specified gate-output variable, computing its gate ID,
locating its up-to-21 occurrences and any corresponding
fixed-width output-bit address costs $O(\log^3G)$ bit operations
and $O(\log G)$ index/work bits on a standard bit-cost
implementation of integer division/multiplication. A physical
input read in a random-access TM also pays to write its binary
address, typically $O(\log T)$ bits when polynomially bounded.
A $\Theta(\log N)$-word RAM comparison must additionally charge
word-address arithmetic and the restrictions on word operations.
No free whole-word semantic SAT primitive is granted.

\paragraph{Fixed-$k$ sequential-input TM $M_k$.}
Let $T_k,S_k,Q_k,C_k$ record physical machine transitions,
visited writable work cells, actual source/work probes and
head-boundary crossing events. If a virtual circuit or child
is evaluated, all source reads, virtual output-bit construction,
failed adaptive choices, retained ancestors, and work-tape head
movements count in this \emph{same} machine. Probes/crossings
are subcounters within $T_k$, not extra additive runtime.
The parameter-only locator does not establish an
$O(\log G)$-space \emph{evaluation} of an arbitrary $G$-gate
DAG, nor a cheap sequential tape path to every indexed input.

\paragraph{One shared read/write tape $M_1$.}
Here the $N$ input cells are part of the footprint $P_1\ge N$;
$S_k$ excludes its read-only source. For fixed $k$, a genuine
single-head sweep simulation conservatively gives
\[
 T_1=O\!\left(T_k(N+S_k+\log(T_k+2))\right),\quad
 P_1=O\!\left(N+S_k+\log(T_k+2)\right).                 \tag{13}
\]
The old independent explicit tape-to-tape audit 171 checked
9,320 matched source/target transition states on a fixed
Boolean comparison verifier; this run does not relabel them as
new measurements. The Hennie--Stearns $O(T_k\log T_k)$ theorem
constructs a \emph{two}-tape simulator, not the M1 result.

\paragraph{Sequential head movement and bounded jumps.}
On the separately specified canonical exact22 easy equality
promise, HEADMOVE-1 establishes the one-sequential-input-head
necessary inequalities
$(R+1)(S+1)=\Omega(N/\log N)$ and
$H(S+1)=\Omega(N^2/\log N)$ with a complete-work-state
crossing/splicing argument, still subject to its declared
source locality and independent review. With a separately
charged input jump radius $J=N^{j+o(1)}$, a sequential
simulation yields the model-specific necessary
$t+s+j\ge2$. Two independently movable input heads or indexed
input can defeat the crossing obstruction on that
\emph{easy} promise. No tape stripe is promoted to FLMV/RAM.

\paragraph{Explicit construction/write floor.}
A Boolean circuit with $G$ distinct two-input gate IDs
requires $\Omega(G\log G)$ physical output bits for the
stated fixed-width gate serialization when written in full.
The proper compiler yields $m_0=13G+4N+12$ clauses and
$n_0=9G+3N+q+8$ names; the previously verified two-level
signed exact-$(2,2)$ regularizer counts
\[
     m_2=4048m_0+5888n_0.                               \tag{14}
\]
Under the pinned architecture, $m_2\ge33856N$ is an
\emph{expansion floor}, not a child of length $N^\lambda$
for $\lambda<1$. This run checks algebraic counts, but does not
materialize or rerun the entire R15\textsuperscript{2} transport.
Its outputs are signed exact-$(2,2)$, not automatically the
stricter clause-monotone class.

\paragraph{Decision versus verifier.}
This construction gives $V_{q,m}(x,z)$, not
$D_N(x)=\bigvee_zV_{q,m}(x,z)$ as a small explicit
all-$N$-bit SAT \emph{decision} circuit. The $q=1.29$ type
of instance-specific boundary relation is likewise not a
full SAT decision family. No all-input source-authenticated
UNSAT certificate cap, unknown-$K$ contraction, uniform ordinary
SAT shortening, ARC-EDGE-90, or $P\ne NP$ follows.

\section{Cross-track handoff and next smallest obligations}
The direct scan/reverse interface fills the final-forward,
scan-output and reverse-output sectors previously marked open
by CIRCUIT-RECORD-OUTPUT-SELECT169. Still missing are the direct
consumers of \emph{all} remaining scan-internal and row-checker
Boolean gates, high-fanout source alias/constant terminals,
complete source input-pin incidence, and inverse Rank/Select
through both proper signed exact-$(2,2)$ regularizer levels.
Even a complete verifier accessor would be only a source
transformation without a charged strictly shorter ordinary-SAT
child and the required deterministic strategy choice from an
\emph{unknown} polynomial-time SAT solver $A$.

The smallest next theorem is arithmetic inverse access for
\emph{primary aliases and source constants}, including the
assignment/$q$-dependent fanout across the full bitonic
network and the final range-checker. A separate falsifier
should test an ordinary variable-width SAT parser with sparse
IDs under one indexed readonly input, explicitly charging
rank/select and malformed encodings.


\clearpage
\section{Run 174 addendum: all primary aliases and all local compiler variable classes}
\noindent\textbf{New finding, fixed architecture and independent referee pending.}
The material before this section is the exact unchanged Run 173 mathematical
predecessor, retained as part of the self-contained source audit. Run 174 now
closes \emph{all nonconstant primary source and witness aliases}, all gate
argument clones, all local gate padding names, and every source/final pin
auxiliary. It additionally reconstructs \emph{every} gate's AND/OR opcode from
its global ID. This does \emph{not} close high-fanout constants or the inverse
consumer map for internal scan and row-checker gate outputs.

\paragraph{Falsifiers proposed before source review.}
The four candidate mechanisms were (H1) scalar scan-internal inverse,
still OPEN; (H2) bounded direct fanout for all primary source aliases,
established here; (H3) arithmetic clone/pad/pin inversion and all-gate opcode
classification, established here; (H4) general arbitrary-source inverse
blocked by the already established adversarial permuted-block example of
SOURCE--SELECT--BARRIER172, not new to this run.

\paragraph{Primary source identification.}
Fix $q\ge3$, $m\ge1$, $w=\lceil\log_2(q+1)\rceil$ and
$N=m[1+3(w+1)]$. Let the global row-checker gate start be
$b_c=G_f+G_s+G_r+c(36w+15)$, as in (2).
For original physical input name $v\in[1,N]$, set
$L=1+3(w+1)$, $c=\lfloor(v-1)/L\rfloor$, $d=(v-1)\bmod L$.
If $d=0$ it is the active bit with \emph{one} direct use,
$(b_c+36w+14,0,-1)$.
Otherwise define $j=\lfloor(d-1)/(w+1)\rfloor$ and
$i=(d-1)\bmod(w+1)$, and put $p=q+3c+j$.
If $i<w$ this is the $i$-th identifier bit of request $j$,
and if $i=w$ this is the request's polarity/payload bit.
Witness source name $v=N+u$, $1\le u\le q$, is payload field $w+2$
of initial assignment record position $p=u-1$.

\paragraph{Exact initial sorter uses.}
At first layer pair rank $r=\lfloor p/2\rfloor$, gate base
$b=rF$, receiving side $\epsilon=p\bmod2$, ascending test
$\mathrm{up}=[(2r)\mathbin{\&}2=0]$, and $D=w+3$, $H=7(w+2)$.
Every identifier bit $0\le i<w$, with $t=w-1-i$, has the three
comparison uses
\[
\begin{cases}
(b+7t,0,-),(b+7t+3,0,+),(b+7t+4,0,-),&\epsilon=\mathrm{up},\\
(b+7t,1,+),(b+7t+3,1,-),(b+7t+4,1,+),&\epsilon\ne\mathrm{up},
\end{cases}
\]
plus the two MUX uses
\[
\begin{cases}
(b+H+3i,1,+),(b+H+3(D+i)+1,1,+),&\epsilon=0,\\
(b+H+3i+1,1,+),(b+H+3(D+i),1,+),&\epsilon=1.
\end{cases}
\]
Payload field $i=w+2$ uses only the MUX pair. These
formulae apply to initial \emph{source} names rather than
previous layer gate outputs: the first layer is their only
direct appearance in the bitonic sorter.

\paragraph{Exact row-checker uses.}
For an original identifier bit at request $(c,j,i)$ the checker
request module begins at $R=b_c+j(8w+3)$.
Its source bit appears once in the initial $w$-input OR
reduction, whose first consumer can be found with a scalar
stream of pair-reduction rounds:
at a round with $n$ live entries and leaf position $i$,
if $i<2\lfloor n/2\rfloor$, use current gate
$\lfloor i/2\rfloor$, slot $i\bmod2$, plus all preceding
round gate counts; if it is the odd unmatched final leaf,
carry it to the next round. This ends in $O(\log w)$ rounds
and stores only three logarithmic-size counters.

Next, the bit enters the $7w$-gate
$\operatorname{lt}(q,\mathrm{id})$ block as its second argument,
at $R+(w-1)+7(w-1-i)+\{0,3,4\}$, operand slot $1$ and
signs $(+,-,+)$. Finally, the two pairwise ID equalities
out of $(0,1),(0,2),(1,2)$ containing source request $j$
each read the bit twice. If $a$ is that pair's index, their
gates are
\[
 b_c+3(8w+3)+a(4w-1)+3i+\{0,1\}.
\]
The first request in the pair occupies operand slot $0$
with signs $(+,-)$; the second occupies slot $1$ with
signs $(-,+)$. This is exactly $1+3+4=8$ checker uses
plus five sorter uses, so each original identifier bit
has \emph{thirteen} direct references.
An original source polarity bit has the two sorter MUX uses
and two checker EQ(sign,lookup) uses at
$R+(w-1)+7w+1$ and $R+(w-1)+7w+2$, operand slot $0$
with signs $(+,-)$, making four. The original active bit
has one, a witness bit two.

\begin{theorem}[Primary variable inverse, new Run 174]
For every $q\ge3,m\ge1$ and every physical primary source
name $v\in[1,N+q]$, a uniform parameter-only procedure
computes exact sorted signed gate-consumer Select and Rank
in $O(\log^3G)$ elementary bit operations and
$O(\log G)$ combined writable/index bits, without incidence
tables. The four direct fanouts by type are $(1,13,4,2)$
for active, identifier, polarity, witness. The proper compiler
source variable has exactly $4f+4$ signed literal occurrences
for an original $x$ bit and $4f$ for an original witness $z$ bit.
The $+4$ is the four explicit original-source pin rows:
$13G+4(v-1)+r$ at literal slot zero, whose sign is determined
by the \emph{actual} source input bit $x_v$.
\end{theorem}
\begin{proof}
The emitter introduces the original source wires only in the
first-layer record and in the row checker's original source
clause fields. Later sorting, scanning, and inverse routing
carry distinct Boolean-gate output variables. The comparator
seven-gate compare and three-gate MUX patterns exhaust its
first-layer reads. In the checker each source identifier
appears once in nonzero OR, three times in less, and twice in
each of two pairwise equality checks, with gate ordering as
displayed. The source sign appears twice in equality;
the active bit in only the final row OR. Every signed gate
operand produces four proper-CNF clone-equality rows at
$13g+4k+\{0,1,2,3\}$, literal slot one, with signs
$(\tau,\tau,-\tau,-\tau)$. Source $x$ bits alone have
four further explicit pin rows. No remaining emitter module
reads a primary source name. At most thirteen signed uses
and three streaming $O(\log G)$-bit counters are live.
Standard schoolbook arithmetic gives the bit-time bound.
\end{proof}

\paragraph{Every gate operation and source-local name.}
Let $s=N+q+3$. Gate $g$ occupies source-variable names
$s+9g+d$ for $d=0,\ldots,8$: argument clones $d=0,1$,
output $d=2$, and padding $d=3,\ldots,8$.
The operation of \emph{every} Boolean gate is arithmetic:
forward comparison pattern AND,AND,OR,AND,AND,OR,AND;
forward and reverse MUX pattern AND,AND,OR;
scan first two kind EQ tests in seven-gate groups,
then for $p>0$ its $3w$ XOR gates and $w-1$ AND tree,
then AND,AND,AND,OR,AND; checker each of three request
blocks is $w-1$ ORs, $7w$ comparator pattern, one
validity AND, AND,AND,OR truth check. Three pairwise
EQ blocks are each $3w$ XOR gates followed by $w-1$ ANDs.
The row ends with five ANDs, two ORs, AND, OR;
the final output tree is all AND. All block boundaries
are the exact arithmetic lengths in (2).

Let $\sigma=-1$ for AND and $\sigma=+1$ for OR.
The own source-CNF occurrences of argument clone $d=0$
at gate $g$ have relative $(\mathrm{row},\mathrm{slot},\mathrm{sign})$
\[
(0,0,-),(1,0,-),(2,0,+),(3,0,+),
(8,0,\sigma),(9,0,-\sigma),(10,0,-\sigma),
\]
and those of clone $d=1$ have
\[
(4,0,-),(5,0,-),(6,0,+),(7,0,+),
(8,1,\sigma),(11,0,-\sigma),(12,0,-\sigma).
\]
Each padding variable $d=3,\ldots,8$ has just one
positive and one negative literal, both in slot $2$,
in relative row pairs $(0,1),(2,3),(4,5),(6,7),
(9,10),(11,12)$ respectively.
Absolute clause indices add $13g$.
After gate variables, $p_0=s+9G$ begins $2N$
source-bit pin auxiliaries followed by six final pin
auxiliaries. Each pin pair has exactly four physical
clauses: the first aux occurs in slot one with signs
$(+,+,-,-)$, the second in slot two with signs
$(+,-,+,-)$. Source pin clause group $i$ begins at
$13G+4i$; final pin group $r\in\{0,1,2\}$ begins at
$13G+4N+4r$.

\begin{theorem}[Compiler-local source inverse, new Run 174]
For all $q\ge3,m\ge1$, every gate's two argument clone
variables and six padding variables, plus every
source-bit and final-output pin auxiliary, has exact
table-free proper-source signed OccSelect and OccRank,
with $O(\log^3G)$ bit time and $O(\log G)$ writable/index
bits. The gate's output variable $d=2$ and original
global constant names are \emph{not} covered by this theorem.
\end{theorem}
\begin{proof}
The source compiler's thirteen rows explicitly list each
private name only in the indicated fixed rows. Gate opcodes
are determined by the mutually disjoint arithmetic sectors
in (2), their per-module offset and the fixed emitter grammar.
The final pin auxiliary ranges are disjoint from all gate-local
and primary ranges. A constant number of logarithmic-size
integers suffices for all rows and signed fields; none of the
$O(G)$ gate descriptors or occurrences is materialized.
\end{proof}

\begin{counterexample}[Wire-to-clause conflation]
For $q=3,m=1$, source active bit $v=1$ has exactly one
direct Boolean use $(1823,0,-)$ but eight physical
proper-CNF occurrences. Original identifier bit $v=2$
has thirteen Boolean gate uses but 56 literal occurrences.
Deleting the source-pin rows incorrectly subtracts
four occurrences per original input bit.
\end{counterexample}

\paragraph{Finite and resource audit.}
Seven exact complete circuits $(q,m)=(3,1),(4,1),(3,2),
(5,2),(7,2),(6,3),(9,2)$ checked 38,528 Boolean gate
opcodes, 1,725 signed primary references, 7,564
physical signed source-input occurrences, and 1,003,224
clone/pad/pin source occurrences against independent
materialized actual compiler rows, all with zero failures.
The earlier 8,192 exhaustive input/witness semantics and
eight-clause UNSAT controls reran identically
(raw SHA256 5eff02d7d9abc9eb2d1b0efb54a15770d5f71cfeca6f2f8a7dd2290fb7c4ccc4).
The prior one-tape simulation reran 340 pairs and 9,320
matched source/simulator states (raw SHA256
3543f3695f2990c4586168ba4cf8cf1e2adecd5d40244f1700bc00419349b80b).

The three native model ledgers remain separated. Fixed-$k$
sequential-input multitape to single read/write tape:
$T_1=O(T_k(N+S_k+\log(T_k+2)))$,
$P_1=O(N+S_k+\log(T_k+2))$, an upper
simulation only; initial input is charged to $P_1$.
Hennie--Stearns is a \emph{two}-tape fast simulation.
The sequential-head exact22 reversal and bounded-jump
bounds do not apply to indexed RAM. The actual source
compiler still expands:
$n_0=9G+3N+q+8$, $m_0=13G+4N+12$,
$m_2=4048m_0+5888n_0$. This is no $N^\lambda$
($\lambda<1$) ordinary-SAT child. Full ordinary-SAT
decision circuitry and unknown-$K$ all-input ARC
resource contraction remain OPEN. P versus NP remains OPEN.

\begin{thebibliography}{9}
\bibitem{flmv} L. Fortnow, R. Lipton, D. van Melkebeek, A. Viglas,
\emph{Time-Space Lower Bounds for Satisfiability},
JACM 52(6), 835--865 (2005),
\url{https://doi.org/10.1145/1101821.1101822}.
\bibitem{pf} N. Pippenger, M. Fischer,
\emph{Relations Among Complexity Measures}, JACM 26, 361--381
(1979), \url{https://research.ibm.com/publications/relations-among-complexity-measures}.
\bibitem{williams} R. R. Williams,
\emph{Simulating Time with Square-Root Space}, STOC 2025,
\url{https://doi.org/10.1145/3717823.3718225}.
\bibitem{hs} F. Hennie, R. Stearns,
\emph{Two-Tape Simulation of Multitape Turing Machines},
JACM 13 (1966), \url{https://doi.org/10.1145/321356.321362}.
\bibitem{prior} Project originals: CIRCUIT-OBLIVIOUS-INVERSE-SORT166,
CIRCUIT-LOCAL-CONSUMER168, CIRCUIT-RECORD-OUTPUT-SELECT169;
roles 1,3,4,5 current reports; DUAL-TM, RAM-VS-TAPE and
HEADMOVE-1 governance directives (2026-10-26).
\end{thebibliography}
\end{document}
% END_OF_VERIFIED_REPORT::LSPV-2::CIRCUIT-PRIMARY-SOURCE-SELECT174
